% GENERATED FILE. Edit canonical lessons, then run npm run generate:books.
% canonical-source-hash: fnv1a64:344b62b5971bcb00
% canonical-lessons: RU-W258-own-line

\chapter{Writing --- My Own Line}
\label{ch:ru-own-line}
\hlchaptermodality{258}

\begin{chapteropening}
\textbf{By the end of this chapter:} \emph{I can write one line about myself by hand from a closed supply of known words, with nothing to copy, and check \ru{л}, \ru{м}, \ru{т} and \ru{я} before anything else.}

Complete RU-W258-own-line: \ru{я}, a verb with its I-ending and the rest, chosen rather than copied, with the handwritten shapes checked first.
\end{chapteropening}

\section[(my own line)]{Your own line — the first writing with no model}

\label{lesson:RU-W258-own-line}



\begin{quote}
Everything you have written so far had something to reach for — a letter to trace, a word to copy, a word you heard.

This does not. You choose the words, from a supply that is closed: only what this book has taught you.
\end{quote}

\subsection*{Writing — controlled composition}
Write \textbf{one line} about yourself, by hand, on paper. Start with \textbf{\ru{я}}, then a verb you own with its \emph{I}-ending, then the rest. End on a dot.

Copying hid one thing from you: the model always had the right shapes in it. Now your hand chooses them, and three are worth checking on purpose.

\textbf{\ru{л}} and \textbf{\ru{м}} — \textbf{\ru{л}} comes to a \textbf{point}, \textbf{\ru{м}} has a \textbf{dip}.

\textbf{\ru{т}} — if your hand writes it like an \emph{m}, give it the bar over the top. The bar is a habit, not a rule, but it keeps the shape apart from \textbf{\ru{ш}} in a run of humps.

\textbf{\ru{я}} — it faces the other way from Latin \emph{R}. A hand that knows English will try to turn it round.

Write the line. Then read it back and check those three letters \textbf{before} you check anything else.

\begin{tcolorbox}[breakable,colback=teal!4,colframe=teal!35!black,title={Before you move on}]
What tells \textbf{\ru{л}} from \textbf{\ru{м}} by hand? (The top — a \textbf{point} or a \textbf{dip}.)

And what did you choose here that the earlier lessons chose for you? (\textbf{The words}, and so \textbf{the shapes}.)
\end{tcolorbox}
